\section{Algebraic singletons and the degree of exact optimizers}
\label{sec:algebraic}

The previous sections ask how hard it is to compare an exact optimal value or
coordinate with a rational threshold. This section asks a prior question:
what kind of number can an exact optimizer be, and how large must an expanded
description of it be? We study the simplest setting in which the answer is not
obvious. The objective $f$ is an unconstrained rational quartic, globally
strongly convex, certified convex by rational data, and with minimum value
zero. Its unique minimizer $p$ is then the unique real zero of $f$, and the
zero sublevel set $\{f\le0\}$ is the singleton $\{p\}$.

The answers are as follows.
\begin{enumerate}[label=(\roman*)]
\item Already in two variables the zero can be irrational, so a single
  strongly convex quartic inequality with integer coefficients can have no
  rational solution (Proposition~\ref{prop:algebraic-bivariate}).
\item The coordinates that occur are exactly the real algebraic numbers with
  one real conjugate. The same class arises for every singleton defined by
  rational convex inequalities, for every unique minimizer of a rational convex
  polynomial, and even for every rational polynomial with a single real zero
  (Theorem~\ref{thm:singleton-field}). A number of degree $d$ is realized in
  $(d+1)/2$ variables, in time polynomial in its dense minimal polynomial; this
  count is optimal for consecutive power coordinates
  (Theorem~\ref{thm:algebraic-dimension}).
\item The degree can be exponential in the number of variables while all
  coefficients have $O(\log n)$ bits (Theorem~\ref{thm:algebraic-cyclic}).
  The construction uses $n+1$ integer squares, the least possible number even
  over $\R$ (Theorem~\ref{thm:sos-length}).
\item For a rational sum of quadratic squares with a single real zero, at
  which the Hessian is positive definite, the field degree of the zero is at
  most $2^n-1$; global convexity improves this to $2^n-3$ for $n\ge3$ and
  $2^n-5$ for $n\ge4$. For $n=5$ the sharp bound is $21$.
  The construction in (iii) is therefore extremal for
  $n\le5$ (Theorem~\ref{thm:algebraic-degree-bounds}).
\item Univariate convex realizations can require exponential degree, although
  an auxiliary variable or a compressed circuit avoids that cost
  (Theorem~\ref{thm:algebraic-univariate}).
\item The optimizer of a certified quartic can be realized together with all
  of its quadratic products, again by a polynomial-size certified quartic
  (Theorem~\ref{thm:quadratic-graph}).
\end{enumerate}
These results matter for exact output. A rational convexity certificate does
not produce rational feasible points, does not bound the algebraic degree of
the optimizer by a constant, and does not make dense minimal-polynomial output
short. They also supply the construction (Lemma~\ref{lem:quartic-realization})
that the lower reductions of Section~\ref{sec:reductions} and the height and
field constructions of Sections~\ref{sec:heights} and~\ref{sec:fields} use to
place a prescribed point at the zero of a certified quartic. None of these
results is a lower bound for exact decision problems; they concern which
exact outputs exist and how long their expanded forms must be. Complete proofs
are in Appendix~\ref{app:algebraic}.

\subsection{Hessian Gram matrices and output formats}
\label{sec:algebraic-models}

In this section a \emph{quadratic} is a polynomial of degree at most two. For
$x,v\in\R^n$ put $w(x,v)=(v,\,x\otimes v)\in\R^{n+n^2}$, where
$x\otimes v$ has entries $x_kv_j$ indexed by ordered pairs $(k,j)$. As in
Section~\ref{sec:heights}, a symmetric matrix $A$ is a \emph{full Hessian
Gram} of a polynomial $f$ of degree at most four if
$v^{\mathsf T}\nabla^2f(x)v=w(x,v)^{\mathsf T}Aw(x,v)$ for all $x,v$. If
$A\succeq\lambda I$ then $\nabla^2f\succeq\lambda I$ everywhere, since
$\norm{w(x,v)}\ge\norm v$. A full Hessian Gram certifies the Hessian; it is
not a Gram matrix of $f$ itself. At a zero of $f$, no polynomial Gram matrix
of $f$ can be positive definite, while a full Hessian Gram can be.

Sums of squares of quadratics carry a canonical full Hessian Gram. For a
quadratic $g(x)=g_0+b^{\mathsf T}x+x^{\mathsf T}Tx$ with $T$ symmetric,
define the $n\times n^2$ matrix $\Delta(b,T)$ and the symmetric matrices
$\Xi(T)$ and $\Gamma(g)$ by
\begin{equation}
\begin{gathered}
 \Delta(b,T)_{i,(k,j)}=2b_kT_{ij}+4b_iT_{kj},\qquad
 \Xi(T)=8\operatorname{vec}(T)\operatorname{vec}(T)^{\mathsf T}+4\,T\otimes T,\\
 \Gamma(g)=\begin{pmatrix}2bb^{\mathsf T}+4g_0T&\Delta(b,T)\\
            \Delta(b,T)^{\mathsf T}&\Xi(T)\end{pmatrix}.
\end{gathered}
 \label{eq:algebraic-canonical}
\end{equation}
Here $\operatorname{vec}(T)_{(k,j)}=T_{kj}$ and
$(T\otimes T)_{(k,j),(k',j')}=T_{kk'}T_{jj'}$. The map $g\mapsto\Gamma(g)$ is
quadratic in the coefficients of $g$; let $\widehat\Gamma(g,g')$ be its
symmetric bilinear polarization. For a polynomial written as
$f=\sum_{k,l}S_{kl}g_kg_l$ with quadratics $g_k$ and a symmetric matrix $S$,
the \emph{canonical Hessian Gram} of this representation is
$\Gamma_S(g)=\sum_{k,l}S_{kl}\widehat\Gamma(g_k,g_l)$. For a weighted sum of
squares $f=\sum_kc_kg_k^2$ it is $\sum_kc_k\Gamma(g_k)$.

\begin{lemma}[Canonical Hessian Grams]
\label{lem:algebraic-gram}
Let $f=\sum_{k,l}S_{kl}g_kg_l$ as above.
\begin{enumerate}[label=(\alph*)]
\item $v^{\mathsf T}\nabla^2f(x)v=w(x,v)^{\mathsf T}\Gamma_S(g)w(x,v)$ for all
  $x,v\in\R^n$.
\item \emph{Covariance.} For $c\in\R^n$ let
  $\Psi_c=\begin{pmatrix}I_n&0\\-c\otimes I_n&I_{n^2}\end{pmatrix}$, where
  $(c\otimes I_n)v=c\otimes v$. Then, as an identity of matrices,
  $\Gamma_S(g)=\Psi_c^{\mathsf T}\,\Gamma_S\bigl(g(\cdot+c)\bigr)\,\Psi_c$.
\item Consequently, if the canonical Gram of the recentred representation
  satisfies $\Gamma_S(g(\cdot+c))\succeq\lambda I$, then
  $\Gamma_S(g)\succeq\lambda(1+\norm c)^{-2}I$.
\end{enumerate}
\end{lemma}

Part (b) is the reason the constructions below need no second approximation.
Positivity is proved for the Gram centred at the common zero $p$ of the square
factors, where they have no constant terms. Covariance then shows that the matrix
$\Gamma_S(g)$, computed directly from the rational square factors, is the
congruent image of that positive definite matrix. It is rational, positive
definite, and of polynomial bit length, and no approximation of $p$ is needed
to obtain it.

\paragraph{Output formats.}
For an algebraic point $p$ we write $\Q(p)=\Q(p_1,\ldots,p_n)$ and call
$[\Q(p):\Q]$ the \emph{field degree} of $p$. A coordinate can be output by the
\emph{dense} coefficient list of its minimal polynomial or by an
\emph{ordinary sparse} list of its nonzero terms
(Definition~\ref{def:models-representations}). The rational circuits of
Definition~\ref{def:models-circuits} compute only rational numbers, so for
irrational coordinates we also use \emph{root circuits}: rational circuits
that may in addition contain gates $z\mapsto z^{1/k}$, returning the positive
real $k$-th root of a positive input, with the integer $k\ge2$ written in
binary. A coordinate of degree $D$ has a dense minimal-polynomial list with
$D+1$ entries. The field degree of $p$ bounds the degree of every coordinate,
but a coordinate need not attain it, and degree alone says nothing about
sparse lists or root circuits. The lower bounds below always name the
coordinate and the format they concern.

\subsection{An integer quartic with an irrational zero}
\label{sec:algebraic-bivariate}

\begin{proposition}[Two variables suffice]
\label{prop:algebraic-bivariate}
Let $q_1=x^2-y$, $q_2=y^2-2x$,
\[
 A=12599x^2-10000xy+7937y^2-15874x-12599y+20000,
 \qquad F=A^2+(100q_1)^2+(100q_2)^2 .
\]
Then $F\in\Z[x,y]$ has coefficients of absolute value less than $2^{30}$, and:
\begin{enumerate}[label=(\alph*)]
\item $\nabla^2F(x,y)\succeq4124\,I$ for all $(x,y)\in\R^2$;
\item $F^{-1}(0)=\{F\le0\}=\{(2^{1/3},2^{2/3})\}$, so $\{F\le0\}$ contains no
  rational point;
\item the canonical Hessian Gram $\Gamma_F$ of the displayed three squares is
  a positive definite integer $6\times6$ matrix, printed in
  \eqref{eq:algebraic-bivariate-gram};
\item the biform $v^{\mathsf T}(\nabla^2F(x)-4096I)v$ has the explicit
  representation \eqref{eq:algebraic-bivariate-squares} as a positive integer
  combination of squares of polynomials with rational coefficients.
\end{enumerate}
\end{proposition}

The construction starts from the three quadratics $q_1,q_2$ and
$q_3=(x-y)^2-2x-y+4$, which vanish at $p=(r,r^2)$, $r=2^{1/3}$. The real
combination $(2r-1)q_1+(r^2-1)q_2+q_3$ equals
$(x-r,y-r^2)\left(\begin{smallmatrix}2r&-1\\-1&r^2\end{smallmatrix}\right)
(x-r,y-r^2)^{\mathsf T}$, a positive definite quadratic form centred at $p$;
the rational combination $A=7599q_1+2937q_2+5000q_3$ approximates $5000$
times it. Thus $A$ vanishes at $p$ exactly, has a positive definite quadratic
part, and has a very small gradient at $p$. Squaring $A$ supplies curvature far
from $p$; the two small squares supply curvature near $p$. Part (a) is proved
by an explicit global estimate, part (c) by
Lemma~\ref{lem:algebraic-gram} and a Schur complement bound, and part (d)
records a direct rational certificate of a slightly weaker modulus.

In version~1 of \cite{SlotSteurerWiedmer2025}, Table~1 leaves open whether
exact feasibility of globally convex rational quartic inequalities always has
a rational witness of polynomial bit length; Appendix~C there gives a convex
sextic with an irrational unique zero. Proposition~\ref{prop:algebraic-bivariate}
answers the quartic witness question negatively in two variables. It says
nothing about the complexity of the decision problem: the point has the short
algebraic description $(2^{1/3},2^{2/3})$.

\subsection{Which algebraic numbers occur}
\label{sec:algebraic-field}

\begin{lemma}[One real conjugate]
\label{lem:one-real-conjugate}
Let $\{p\}\subseteq\R^n$ be a singleton that is either the real zero set of
finitely many polynomials with rational coefficients, or the set
$\{x:f_i(x)\le0,\ 1\le i\le s\}$ for finitely many rational convex polynomials
$f_i$. Then $p$ has real algebraic coordinates, the field $K=\Q(p)$ has exactly
one embedding into $\R$, $[K:\Q]$ is odd, and the minimal polynomial of every
coordinate $p_j$ has exactly one real root.
\end{lemma}

For the zero-set case, every real embedding of $K$ maps $p$ to another real
common zero, hence fixes $p$. In the convex case the inequalities active at
$p$ already define $\{p\}$, and the same argument applies to them. The passage
from the joint field to a single coordinate uses that $[K:\Q(p_j)]$ is odd, so
every real embedding of $\Q(p_j)$ extends to a real embedding of $K$.

\begin{theorem}[Coordinates of rational convex singletons]
\label{thm:singleton-field}
For a real number $\alpha$ the following are equivalent.
\begin{enumerate}[label=(\roman*)]
\item $\alpha$ is algebraic and its minimal polynomial over $\Q$ has exactly
  one real root.
\item For some $n$, some $f\in\Q[x_1,\ldots,x_n]$ has a single real zero, and
  $\alpha$ is one of its coordinates.
\item For some $n$, finitely many rational convex polynomials $f_i$ on $\R^n$
  satisfy $\{x:f_i(x)\le0\ \forall i\}=\{p\}$ with $\alpha$ a coordinate
  of $p$.
\item For some $n$, $\alpha$ is a coordinate of the unique minimizer of a
  rational convex polynomial on $\R^n$ that has exactly one minimizer.
\item For some $n$, $\alpha$ is a coordinate of the unique zero of a rational
  quartic $f$ on $\R^n$ that is a sum of $n+1$ squares of rational quadratics,
  satisfies $\nabla^2f\succeq I$ on $\R^n$, and has a positive definite
  rational full Hessian Gram.
\end{enumerate}
If (i) holds and $\alpha$ has degree $d\ge3$, then (v) holds with
$n=(d+1)/2$ and zero $(\alpha,\alpha^2,\ldots,\alpha^n)$, and $f$, its square
factors, and its Gram matrix are computed from the dense coefficient list of
any rational multiple of the minimal polynomial in time polynomial in that
list's binary length. If $\alpha$ is rational, (v) holds with $n=1$ and
$f=(x-\alpha)^2+(x-\alpha)^4$.
\end{theorem}

The theorem says that convexity, certified convexity, degree four, and a sum
of squares representation cost nothing in the class of reachable coordinates:
they reach exactly the numbers that any rational polynomial with a single real
zero can reach. The restriction to one real conjugate concerns singletons and
unique unconstrained minimizers. It does not apply to constrained optimizers:
$-x$ has the unique minimizer $\sqrt2$ on the rational convex set
$\{x:x^2\le2\}$, and $\sqrt2$ has two real conjugates. In particular, no
singleton defined by rational convex inequalities has $\sqrt2$ as a
coordinate, which is why Square Root Sum comparisons are not encoded by such
singletons (Section~\ref{sec:reductions}).

Several convex quadratic inequalities suffice when one prefers quadratics to
one quartic. The following corollary realizes power coordinates of degree $d$
with three rational ellipsoids, so the Hessian span of a convex quadratic
system with an irrational singleton can be three while the degree of the
singleton is unbounded.

\begin{corollary}[Three rational ellipsoids]
\label{cor:algebraic-three-quadrics}
Let $P\in\Q[T]$ be irreducible of degree $d\ge3$ with exactly one real root
$\alpha$, given by its dense coefficient list. In time polynomial in the
length of that list one can compute rational quadratics $g_1,g_2,g_3$ in $d-1$
variables, each with positive definite Hessian, such that
$\{x:g_i(x)\le0,\ i=1,2,3\}=\{(\alpha,\alpha^2,\ldots,\alpha^{d-1})\}$. Each
set $\{g_i\le0\}$ is a full-dimensional ellipsoid.
\end{corollary}

\subsection{A generic quantitative construction}
\label{sec:algebraic-realization}

The realization results of this section, and those of
Sections~\ref{sec:reductions}, \ref{sec:heights} and~\ref{sec:fields}, use
one assembly step. A point $p$ is
described by rational quadratics $r_j$ vanishing at it with a well-conditioned
Jacobian, and by one rational quadratic $g$ vanishing at $p$ whose quadratic
part is positive definite and whose gradient at $p$ is small. Such a $g$
usually comes from an irrational positive semidefinite \emph{exposing}
quadratic that vanishes to second order at $p$; approximating its coefficients
inside the rational space of quadratics that vanish at $p$ keeps $g(p)=0$
exactly and makes $\nabla g(p)$ as small as required.

\begin{lemma}[Quantitative quartic realization]
\label{lem:quartic-realization}
Let $p\in\R^n$, and let $g,r_1,\ldots,r_m\in\Q[x_1,\ldots,x_n]$ be
polynomials of degree at most two that vanish at $p$. Write
\[
 g(p+u)=\ell^{\mathsf T}u+u^{\mathsf T}Hu,\qquad
 r_j(p+u)=b_j^{\mathsf T}u+u^{\mathsf T}T_ju,
\]
with symmetric $H,T_j$. Let $\mu,\Lambda,\beta,\nu$ be positive rationals with
\begin{equation}
 H\succeq\mu I,\quad \norm H\le\Lambda,\quad \norm{T_j}\le1,\quad
 \norm{b_j}\le\beta\ (1\le j\le m),\quad
 \sum_{j=1}^mb_jb_j^{\mathsf T}\succeq\nu^2I .
 \label{eq:algebraic-realization-data}
\end{equation}
Let $t>0$ be rational, put $\varepsilon=t^2$, and suppose
\begin{equation}
 \varepsilon\le\min\Bigl\{1,\ \frac{\mu^2}{2m},\
   \frac{\nu^2\mu^2}{36\,n(\Lambda+m\beta)^2}\Bigr\},
 \qquad \norm\ell\le\varepsilon .
 \label{eq:algebraic-realization-eps}
\end{equation}
Then $f=(g/(t\nu))^2+\sum_{j=1}^m(r_j/\nu)^2$ has the following properties.
\begin{enumerate}[label=(\alph*)]
\item $f\in\Q[x]$ has degree exactly four, is a sum of $m+1$ squares of
  rational quadratics vanishing at $p$, and its quartic form is positive
  definite.
\item $\nabla^2f(x)\succeq\tfrac32I$ for every $x\in\R^n$. Hence $p$ is the
  unique zero and the unique minimizer of $f$.
\item The canonical Hessian Gram $\Gamma_f$ of the displayed representation is
  a rational full Hessian Gram of $f$, and $\Gamma_f\succeq\gamma I$ with
  \[
   \gamma=\frac{\min\{3/2,\ 2\mu^2/(\varepsilon\nu^2)\}}
          {\bigl((1+\theta)^2+1\bigr)(1+\norm p)^2},\qquad
   \theta=\frac{3\sqrt n\,\varepsilon(\Lambda+m\beta)}{\mu^2}.
  \]
  Its entries are fixed polynomial expressions in the coefficients of
  $g,r_1,\ldots,r_m$ and in $t,\nu$; computing it uses
  $O(m(n+n^2)^2)$ rational operations on numbers of the input bit length.
\end{enumerate}
\end{lemma}

The factor $n$ in the last bound of \eqref{eq:algebraic-realization-eps} is
used only for part (c); part (b) holds without it
(Remark~\ref{rem:algebraic-curvature}). The proof computes the
canonical Gram centred at $p$, where it has three blocks: a constant block at
least $2\varepsilon\nu^2I$, a quartic block at least $2\mu^2I$ even though the
individual $T_j\otimes T_j$ can be indefinite, and a small cross block. A
Schur complement estimate gives positivity, and covariance transfers it to the
rational Gram. Squaring an indefinite quadratic does not produce a convex
function, and the estimate retains the corresponding negative terms. Positive
definite Hessian Gram matrices and Schur-complement arguments for SOS-convex
regularization are established techniques \cite{AhmadiChaudhryZhang2024}; the
additional feature here is that a prescribed, possibly irrational, zero is kept
exactly while all coefficients stay rational and of polynomial size. The
hypotheses force $p$ to be algebraic: $n$ of the residuals have a nonsingular
Jacobian at $p$, so a small rational box isolates $p$ among their real common
zeros, and the transfer argument in the proof of
Lemma~\ref{lem:one-real-conjugate} applies. The construction, however, uses
no algebraic representation of $p$, and $p$ need not be known or rational.
The constructions in Sections~\ref{sec:reductions}, \ref{sec:heights},
and~\ref{sec:fields} apply the lemma at the gate values of certified odd-root
circuits and of rational unit-quaternion circuits, at points given by
cube-root chains and by repeated complex squaring, and at points with
coordinates in towers of radicals.

\subsection{Power coordinates and the dimension count}
\label{sec:algebraic-power}

\begin{theorem}[Realization in power coordinates]
\label{thm:algebraic-dimension}
Let $P\in\Q[T]$ be irreducible of degree $d\ge3$ with exactly one real root
$\alpha$, given by its dense list of rational coefficients of total binary
length $L$, and let $n$ be an integer with $(d+1)/2\le n\le d-1$. A
deterministic algorithm running in time polynomial in $L$ computes rational
quadratics $g,r_1,\ldots,r_n$ in $n$ variables and positive rationals
$t,\mu,\Lambda,\beta,\nu$ that satisfy the hypotheses of
Lemma~\ref{lem:quartic-realization} at $p=(\alpha,\alpha^2,\ldots,\alpha^n)$. Hence
$F_P=(g/(t\nu))^2+\sum_{j=1}^n(r_j/\nu)^2$ is a sum of $n+1$ squares of
rational quadratics, $\nabla^2F_P\succeq\tfrac32I$, $F_P^{-1}(0)=\{p\}$, and
the canonical Hessian Gram of $F_P$ is a positive definite rational matrix; all
outputs have bit length polynomial in $L$.

Conversely, if $f\in\Q[x_1,\ldots,x_k]$ has degree at most four, is
nonnegative on $\R^k$, and has the unique real zero
$(\alpha,\alpha^2,\ldots,\alpha^k)$, then $k\ge(d+1)/2$.
\end{theorem}

The algorithm first replaces $P$ by the monic minimal polynomial $P_1$ of
$\alpha$; the input may be any rational multiple of it, including a negative
one. The exposing quadratic comes from the nonnegative univariate polynomial
$(T-\alpha)P_1(T)(1+T^2)^e$, where $e=n-(d+1)/2$ raises the degree to $2n$:
its positive factor $P_1(T)(1+T^2)^e/(T-\alpha)$ is a product of real
quadratics with positive definite $2\times2$ Gram matrices, which a recursion
compresses onto consecutive monomials. Multiplying by the difference matrix
for $T-\alpha$ produces a positive semidefinite matrix $Q_*$ of order $n+1$
whose kernel is spanned by $(1,\alpha,\ldots,\alpha^n)$. Its entries involve
all complex roots, so the algorithm approximates the roots with certified
precision, forms a rational approximation, and projects it exactly onto the
rational space of matrices that represent multiples of $P$. Root separation
from the integer discriminant keeps every required precision at polynomially
many bits. The residuals are $x_1x_j-x_{j+1}$ and one quadratic version of
$P_1(x_1)=0$, whose Jacobian determinant at $p$ is a positive multiple of
$P_1'(\alpha)$.

The lower bound concerns the prescribed consecutive-power point only. A degree
$d$ number has other coordinate encodings, and we do not know whether some of
them need fewer variables. Irreducibility and the single real root are
promises; the algorithm does not test them.

\subsection{Exponential degree with small coefficients}
\label{sec:algebraic-cyclic}

Theorem~\ref{thm:algebraic-dimension} uses about $d/2$ variables for degree
$d$. Far fewer variables suffice for special numbers. Let $n\ge2$, $m=n+1$,
$\eta=(-1)^m$, and
\[
 d_n=\frac{2^m-\eta}{3}=\frac{2^{n+1}-(-1)^{n+1}}3,\qquad
 e_i=\frac{(-2)^i-1}{3}\ (0\le i\le n),\qquad
 p_i=2^{e_i/d_n}.
\]
The sequence $d_n$ begins $3,5,11,21,43,85$. With $x_0=1$ and indices modulo
$m$, the point $p=(p_1,\ldots,p_n)$ is the unique real common zero of the $m$
cyclic binomials $x_i^2-2^{\theta_i}x_{i+1}x_{i+2}$, where $\theta_i=0$ except
$\theta_{m-2}=\eta$ and $\theta_{m-1}=-\eta$. The weights $p_i^{-2}$ combine
these binomials into the anchored cycle energy
$\tfrac12\sum_i(x_i/p_i-x_{i+1}/p_{i+1})^2$, an exposing quadratic whose
curvature and Jacobian are controlled by inverse-polynomial quantities.

\begin{theorem}[Small integer quartics with exponentially large field degree]
\label{thm:algebraic-cyclic}
For every $n\ge2$ a deterministic algorithm running in time polynomial in $n$
computes integer quadratics $s_0,\ldots,s_n$ in $n$ variables, with
coefficients of $O(\log(n+1))$ bits, such that $F_n=\sum_{i=0}^ns_i^2$
satisfies:
\begin{enumerate}[label=(\alph*)]
\item $\nabla^2F_n\succeq I$ on $\R^n$, and the canonical Hessian Gram of
  $\sum_is_i^2$ is a positive definite integer matrix;
\item the unique zero of $F_n$ is $p\in(1/2,2)^n$, and
  $[\Q(p):\Q]=[\Q(p_1):\Q]=d_n$;
\item $F_n$ has $O(n^2)$ monomials, each with an $O(\log(n+1))$-bit
  coefficient;
\item $F_n$ is not a sum of fewer than $n+1$ squares of real polynomials;
\item the condition number of $\nabla^2F_n(p)$ is at most $65600\,n^2$;
\item after the substitution $x_1\mapsto x_1-1$, the first coordinate of the
  unique zero is $1+2^{-1/d_n}$, whose primitive minimal polynomial
  $2(T-1)^{d_n}-1$ has all $d_n+1$ coefficients nonzero and total binary
  length $\Theta(d_n^2)$.
\end{enumerate}
\end{theorem}

The field degree is exponential in the dimension, while the whole instance is
computed in time polynomial in $n$. Even with dense exponent vectors and the
full canonical Hessian Gram of order $n+n^2$ included, the certified input has
binary length $L=O(n^4\log n)$. The dense and ordinary sparse
minimal-polynomial lists in part (f) therefore have length
$2^{\Omega(n)}=2^{\Omega((L/\log L)^{1/4})}$, which is superpolynomial, though
not exponential, in $L$. No algorithm that writes such a list for this
coordinate runs in time polynomial in $L$. Short output in other formats is
not excluded. Each coordinate $p_i=2^{e_i/d_n}$ is determined by its rational
exponent $e_i/d_n$, of $O(n)$ bits. With $a=2^{1/d_n}$, the relation
$e_{i+1}=-2e_i-1$ gives $p_1=1/a$ and $p_{i+1}=1/(ap_i^2)$, so one root gate
with the $O(n)$-bit index $d_n$, applied to $2$, followed by $O(n)$
multiplications and divisions, is a root circuit for the whole optimizer.
Part (e) shows that the large degree does not come from ill conditioning:
the Hessian at the optimizer has condition number polynomial in $n$. This is
a statement at the
optimizer, not a uniform smoothness bound on a neighbourhood.
Part (f) needs the translation, because $p_1=2^{-1/d_n}$ itself has the
two-term minimal polynomial $2T^{d_n}-1$. The grounded determinant $d_n$ of the
directed square cycle and its exponential growth are classical; see
\cite{LoncParolWojciechowski2001}. The supporting construction of the proof
also gives $n+1$ rational ellipsoids with the same singleton intersection.

The number of squares in part (d) is optimal for a general reason.

\begin{theorem}[Minimum number of squares]
\label{thm:sos-length}
Let $f\in\R[x_1,\ldots,x_n]$ be globally convex of degree exactly four, and
suppose $f=\sum_{i=1}^mg_i^2$ with $g_i\in\R[x]$, $f(p)=0$, and
$\nabla^2f(p)\succ0$ for some $p\in\R^n$. Then $m\ge n+1$. The bound is
attained for every $n$ by $(\sum_ix_i^2)^2+\sum_ix_i^2$, by every quartic of
Lemma~\ref{lem:quartic-realization} with $m=n$, and by the quartics $F_n$ of
Theorem~\ref{thm:algebraic-cyclic}.
\end{theorem}

Each hypothesis is needed. The quartic $(1+\norm x^2)^2$ is globally strongly
convex with one square and minimum one; $\norm x^4$ has a degenerate zero and
one square; the nonconvex $x^2+(y-x^2)^2$ has two squares in two variables and
Hessian $2I$ at its unique zero; a positive definite quadratic is a sum of $n$
squares; and the sextic $(x_1+x_1^3)^2+\sum_{j\ge2}x_j^2$ is globally strongly
convex with $n$ squares. The proof shows first that the leading quadratic part
of the square factors does not depend on directions in which the quartic form
vanishes, and that convexity forbids quadratic coupling to those directions;
minimizing over them leaves a square quadratic map with one regular zero and a
proper leading part. A proper even map has even topological degree, whereas a
map with one regular zero has degree $\pm1$.

The cyclic network is optimal among the binomial systems that its exposing
strategy can use.

\begin{proposition}[Optimality within binomial exposing networks]
\label{prop:algebraic-network}
Let $n\ge2$, $m=n+1$, and let $\mathcal D$ be a strongly connected directed
multigraph on $\{0,\ldots,n\}$, loops and parallel arcs allowed, with exactly
two outgoing arcs $i\to a(i)$, $i\to b(i)$ at every vertex. Fix $x_0=1$ and
nonzero rationals $c_i$, and consider the $m$ equations
$x_i^2=c_ix_{a(i)}x_{b(i)}$, $0\le i\le n$. If the system has exactly one real
solution $p$, then $p$ is algebraic and $[\Q(p):\Q]\le d_n$; if $\mathcal D$
is not Eulerian, then $[\Q(p):\Q]\le2^{n-1}-1$. Conversely, if $m$ rational
homogeneous quadratic binomials in $(X_0,\ldots,X_n)$ vanish at a point
$(1,p_1,\ldots,p_n)$ with nonzero coordinates and some real linear
combination of them is a positive semidefinite quadratic form with kernel
spanned by that point, then after reordering and rational rescaling the
binomials have the above form; in particular the field degree is at most $d_n$,
and fewer than $m$ binomials cannot have such a combination.
\end{proposition}

The proof counts complex solutions as characters of the finite group
$\Z^n/\Lambda$, where $\Lambda$ is the row lattice of the exponent matrix;
its order is the gcd $g$ of the rooted spanning-tree counts. Uniqueness of the
real solution is equivalent to $g$ being odd. For Eulerian graphs, classical
identities for the interlace polynomial of an Euler tour give
$3g\le2^m-(-1)^m$ \cite{ArratiaBollobasSorkin2004,BalisterBollobasCutlerPebody2002}.
The description of binomial solution sets by characters of a lattice quotient
is standard \cite{EisenbudSturmfels1996}. The proposition restricts the method,
not the problem: it does not bound the field degree of arbitrary rational sums
of quadratic squares.

\subsection{Upper bounds on the degree}
\label{sec:algebraic-bounds}

How large can the field degree of the zero be? For quartics given as rational
sums of quadratic squares, intersection theory bounds it in terms of $n$. After
a rational change of coordinates and elimination of flat directions, global
convexity leaves the square factors
without common real zeros at infinity, which would otherwise spoil classical
Cayley--Bacharach arguments.

\begin{theorem}[Degree bounds for rational sums of quadratic squares]
\label{thm:algebraic-degree-bounds}
Let $f=\sum_{j=1}^mr_j^2$ with $r_j\in\Q[x_1,\ldots,x_n]$ of degree at most
two, suppose that $f$ has exactly one real zero $p$ and that
$\nabla^2f(p)\succ0$, and put $D=[\Q(p):\Q]$. Then $p$ is algebraic, $D$ is
odd, and:
\begin{enumerate}[label=(\alph*)]
\item $D\le2^n-1$;
\item if $f$ is globally convex and $n\ge3$, then $D\le2^n-3$; if moreover
  $n\ge4$, then $D\le2^n-5$;
\item if $f$ is globally convex and $n=5$, then $D\le21$.
\end{enumerate}
Consequently, among globally convex quartics in this class the largest field
degrees in dimensions $1,2,3,4,5$ are exactly $1,3,5,11,21$; for $n\ge2$ they
are attained by the quartics $F_n$ of Theorem~\ref{thm:algebraic-cyclic}.
\end{theorem}

\begin{table}[ht]
\centering
\begin{tabular}{@{}lccccc@{}}
\toprule
$n$ & 2 & 3 & 4 & 5 & 6\\
\midrule
cyclic degree $d_n$ & 3 & 5 & 11 & 21 & 43\\
upper bound & 3 & 5 & 11 & 21 & 59\\
\bottomrule
\end{tabular}
\caption{Field degrees of zeros of globally convex rational sums of quadratic
squares with a nondegenerate zero. The maximum is known exactly for $n\le5$.
For $n\ge6$ the cyclic family gives $d_n=(2^{n+1}-(-1)^{n+1})/3$ and the upper
bound is $2^n-5$.}
\label{tab:algebraic-degrees}
\end{table}

Part (a) selects $n$ square factors with a nonsingular Jacobian at $p$; all
conjugates of $p$ are simple isolated zeros of these $n$ quadrics, so a
weighted form of B\'ezout's theorem bounds their number by $2^n$, and oddness
gives $2^n-1$. For (b), global convexity lets us remove, by a rational change
of coordinates that preserves the field and the rational squares, every
direction in which the quartic form vanishes. Afterwards the square factors
have no common real zero at infinity, although common complex zeros at
infinity may remain. If the degree were $2^n-1$ or $2^n-3$, the conjugates of
$p$ would leave a residual of length one or three in a complete intersection
of $n$ rational linear combinations of the square factors, chosen generically
for length three, and Cayley--Bacharach theorems
\cite{EisenbudGreenHarris1996} would make that residual contain a second real
common zero. In five variables, the remaining residual
lengths $1,3,5,7,9$ are excluded by a parity argument for a Gorenstein pairing
on the residual scheme, and positive-dimensional complex base components are
excluded by a residual-intersection count over a list of low-degree real-point
free curves and surfaces. These five-variable steps import precisely stated
results of intersection theory \cite{Fulton1998,EklundJostPeterson2013} and
the classification of curves of minimal degree \cite{EisenbudHarris1987}; the
contracts are recorded in Appendix~\ref{app:algebraic-imports}.

Rational square factors are essential to this argument. Without them, the
gradient equations give the following weaker bound, and in two variables a
classification of rational nonnegative ternary quartics recovers the sharp
bound.

\begin{proposition}[Bounds without square factors]
\label{prop:algebraic-node-bound}
Let $f\in\Q[x_1,\ldots,x_n]$ have degree at most four, be nonnegative on
$\R^n$, and have exactly one real zero $p$, with $\nabla^2f(p)$ nonsingular.
Then $p$ is algebraic and $[\Q(p):\Q]$ is odd and at most $2\cdot3^{n-1}-1$.
If $n=2$ and $f$ is globally convex, then $[\Q(p):\Q]\le3$.
\end{proposition}

\begin{lemma}[Rational descent in two variables]
\label{lem:algebraic-bivariate-descent}
Let $f\in\Q[x_1,x_2]$ have degree four, be nonnegative on $\R^2$, have exactly
one real zero, and have a positive definite quartic form. Then $f$ is a sum of
squares of polynomials with rational coefficients. In particular this holds
for every certified quartic in two variables with minimum zero.
\end{lemma}

The lemma uses Scheiderer's description of rational nonnegative ternary
quartics that are not sums of rational squares \cite{Scheiderer2016}: such a
form is a product of four complex lines in general position, and then it has
at least two real projective zeros. The lemma is the reason why the
three-variable rational descent obstruction of Section~\ref{sec:fields} cannot
occur in two variables. Global convexity cannot be dropped from
Theorem~\ref{thm:algebraic-degree-bounds}(b): Example~\ref{ex:algebraic-seven}
in the appendix gives three rational quadratics in three variables whose sum of
squares has a unique real zero of field degree $7>2^3-3$, with positive
definite Hessian there.

\subsection{Univariate realizations: degree, dimension, and format}
\label{sec:algebraic-univariate}

A univariate rational polynomial can also have a single real zero at a
number with one real conjugate; the square of its minimal polynomial does.
Requiring convexity changes the cost.

\begin{theorem}[Univariate convex realizations]
\label{thm:algebraic-univariate}
\begin{enumerate}[label=(\alph*)]
\item A real number is the unique real zero of a rational globally strongly
  convex univariate polynomial if and only if it is algebraic with exactly one
  real conjugate.
\item Let $P\in\Z[T]$ be irreducible of degree $d\ge1$, with positive leading
  coefficient, coefficients of absolute value at most $2^\tau$, and exactly one
  real root $\alpha$. In time polynomial in $d+\tau$ one can compute
  $c\in\Q$, $\kappa\in\Q_{>0}$, and an integer $N\ge2$, each of binary length
  $O(d^3(\tau+\log(d+1)))$, such that $G=\kappa P^2(1+(T-c)^2)^N$ satisfies
  $G''\ge1$ on $\R$ and $G^{-1}(0)=\{\alpha\}$. Thus $G$ is computed by a
  rational arithmetic circuit with input $T$ and size polynomial in $d+\tau$,
  while its degree $2d+2N$
  may be exponential and its value at the integer $2^{\tau+2}$ is at least
  $2^N$.
\item For $k\ge1$ let $P_k=T^3-3T+2+2^{-2k}$. It is irreducible, has a single
  real root $\alpha_k$, and has binary length $O(k)$. Every rational globally
  convex univariate $F$ whose only real zero is $\alpha_k$ satisfies
  $(\deg F)^2>\tfrac342^k$. Every finite system of rational convex univariate
  inequalities with solution set $\{\alpha_k\}$ contains a polynomial of such
  degree. Every nonconstant rational globally convex $J$ with a minimizer at
  $\alpha_k$ satisfies $(\deg J-1)^2>\tfrac342^k$.
\item The two-variable quartic of Theorem~\ref{thm:algebraic-dimension}
  ($d=3$, $n=2$) has unique zero $(\alpha_k,\alpha_k^2)$ and binary length
  polynomial in $k$. The quartic $J_k=T^4/4-3T^2/2+(2+2^{-2k})T$ is globally
  quasiconvex, not convex, and has unique minimizer $\alpha_k$.
\end{enumerate}
\end{theorem}

Parts (a) and (b) multiply $P^2$ by $(1+(T-c)^2)^N$; the exponent $N$ must be chosen
before the rational centre $c$, which must then lie very close to $\alpha$
(a fixed centre does not work in general). Convexification of positive
polynomials by such multipliers is established \cite{KurdykaSpodzieja2015};
the zero of $P^2$ is what forces the moving rational centre. Part (c) uses that
two complex
conjugates of $\alpha_k$ lie within $2^{-k}$ of the real point $1$, at distance
more than three from $\alpha_k$: a convex polynomial vanishing at $\alpha_k$
must vanish at these conjugates, and Markov's inequality on the interval
from $\alpha_k$ to their common real part forces large degree
\cite{BorweinErdelyi1995}. Thus, in one variable, exact
convex realization has an unavoidable exponential degree for some
$O(k)$-bit inputs; one auxiliary variable, a compressed circuit, or
quasiconvexity removes this cost. Part (b) supplies no cheap exact evaluation:
expanded rational values of $G$ at short rational arguments can be
exponentially long.

\subsection{Realizing the quadratic products of an optimizer}
\label{sec:algebraic-graph}

Some constructions need a certified quartic whose zero records not only an
optimizer $p$ but also all products $p_ip_j$. The next theorem provides one
without computing $p$ or any minimal polynomial. It is used later in the
paper to make both blocks of a block-separated certificate obstruction
strongly SOS-convex.

\begin{theorem}[Quadratic graph realization]
\label{thm:quadratic-graph}
Let $f\in\Q[x_1,\ldots,x_n]$ have degree at most four and a rational full
Hessian Gram $A\succ0$, and suppose that $\min f=0$; this is a promise and is
not checked. Let $p$ be the minimizer, $N=n+n(n+1)/2$, and use variables
$(y,z)$ with $z=(z_{ij})_{1\le i\le j\le n}$. In deterministic time polynomial
in the binary length $L$ of $(f,A)$ one can compute rational quadratics
$g,r_1,\ldots,r_N$ in $(y,z)$ and positive rationals $t,\mu,\Lambda,\beta,\nu$
that satisfy the hypotheses of Lemma~\ref{lem:quartic-realization} at
$w_*=(p,(p_ip_j)_{i\le j})$. Hence $B=(g/(t\nu))^2+\sum_{j=1}^N(r_j/\nu)^2$ is
a sum of $N+1$ rational squares, $B^{-1}(0)=\{w_*\}$,
$\nabla^2B\succeq\tfrac32I$, and the canonical Hessian Gram of $B$ is a
positive definite rational matrix. No representation of $B$ uses fewer than
$N+1$ real squares.
\end{theorem}

The exposing quadratic is a quadratic lift of $f$ itself. For a rational
centre $\hat p$, the Taylor formula with integral remainder and the full
Hessian Gram express $f(y)$ as a rational quadratic $g_{\hat p}(y,z)$
evaluated on the graph $z_{ij}=y_iy_j$. Hence $g_{\hat p}(w_*)=f(p)=0$
exactly for every $\hat p$, which is what the assembly lemma needs. Its
quadratic part is uniformly positive definite, and its gradient at $w_*$ is
at most a constant
multiple of $\norm{\hat p-p}$. A rational
$\hat p$ close to $p$ is computed by approximate convex minimization
(Lemma~\ref{lem:convex-value}) and strong convexity. The residuals are the
graph relations $y_iy_j-z_{ij}$ and a quadratic version of $\nabla f(y)$, with
Jacobian determinant $\det\nabla^2f(p)$ at $w_*$.

\subsection{Scope and open questions}
\label{sec:algebraic-scope}

The results of this section are statements about exact outputs and exact
certificates. They give no lower bound for deciding feasibility, comparing
values, or approximating the optimizer. In the explicit families above the
optimizer also has a short description, as a radical expression, a root of
the input polynomial, or the zero of the input quartic, and the comparison
problems for certified quartics are covered by the upper bounds of
Section~\ref{sec:upper}. Dimension-exponential
lower bounds on output length are superpolynomial in the input length of the
families, not exponential in it.

Several questions remain open. The largest field degree of the zero of a
globally convex rational sum of quadratic squares is unknown for $n\ge6$;
the cyclic family gives about $\tfrac232^n$ and
Theorem~\ref{thm:algebraic-degree-bounds} gives $2^n-5$. The sharp bound for
rational convex quartics without a rational square representation is unknown
for $n\ge3$; Proposition~\ref{prop:algebraic-node-bound} gives
$2\cdot3^{n-1}-1$. We do not know the fewest variables needed for a degree $d$
number in coordinate encodings other than consecutive powers, the optimal
degree in Theorem~\ref{thm:algebraic-univariate}(c), or the shortest sparse
or circuit description of the univariate realizations.
